% --- 1. ОБЩИЕ НАСТРОЙКИ ТЕКСТА ---
\onehalfspacing
\setlength{\parindent}{1.25cm} 

% Борьба с дырками в строках и переносами
\sloppy 
\hyphenpenalty=1000 
\tolerance=2000

% --- 2. ПРИНУДИТЕЛЬНОЕ "СОДЕРЖАНИЕ" (Название) ---
\usepackage{etoolbox}
\addto\captionsrussian{%
  \renewcommand{\contentsname}{СОДЕРЖАНИЕ}%
}

% --- 3. ШРИФТЫ ЗАГОЛОВКОВ (Везде 14pt жирный) ---
\setkomafont{disposition}{\rmfamily\bfseries\normalsize}
\setkomafont{chapter}{\rmfamily\bfseries\normalsize}
\setkomafont{section}{\rmfamily\bfseries\normalsize}
\setkomafont{subsection}{\rmfamily\bfseries\normalsize}

% Пункты в оглавлении — обычные (НЕ жирные)
\setkomafont{chapterentry}{\normalfont\normalsize}
\setkomafont{chapterentrypagenumber}{\normalfont\normalsize}
\KOMAoptions{toc=chapterentrydotfill, numbers=noendperiod}

% --- 4. УМНАЯ ЛОГИКА ВЫРАВНИВАНИЯ (Центр для Содержания, Отступ для остального) ---
\makeatletter
% Логика для ГЛАВ (chapter)
\renewcommand{\chapterlinesformat}[3]{%
  \ifstr{#2}{}{% ЕСЛИ НЕТ НОМЕРА (как в Содержании, Введении, Литературе)
    \begin{center}#3\end{center}% ЦЕНТРИРУЕМ
  }{% ЕСЛИ ЕСТЬ НОМЕР (как 1. Обзор технологий)
    \hspace{1.25cm}#2#3% ДЕЛАЕМ ОТСТУП 1.25см
  }%
}

% Логика для РАЗДЕЛОВ (section, subsection)
\renewcommand{\sectionlinesformat}[4]{%
    \hspace{1.25cm}#3#4% Всегда отступ 1.25см
}
\makeatother

% Интервалы между заголовками и текстом
\RedeclareSectionCommand[beforeskip=12pt, afterskip=6pt]{chapter}
\RedeclareSectionCommand[beforeskip=12pt, afterskip=6pt]{section}
\RedeclareSectionCommand[beforeskip=12pt, afterskip=6pt]{subsection}

% --- 5. СПИСКИ (enumitem) ---
\setlist{
	leftmargin=0pt,
	nosep,
	wide=\parindent,
	labelsep=0.5em,
	listparindent=\parindent
}
\setlist[itemize,1]{label=--} 
\setlist[enumerate,1]{label=\arabic*.}

% --- 6. ПОДПИСИ И ТАБЛИЦЫ ---
\captionsetup{justification=centering, singlelinecheck=false}
\DeclareCaptionLabelFormat{simple}{#1~#2}
\captionsetup{labelformat=simple}
\DeclareCaptionLabelSeparator{emdash}{ --- }
\captionsetup[table]{
	labelsep=emdash,
	justification=raggedright,
	singlelinecheck=false,
	position=top,
	skip=6pt
}

% --- 7. ЛИСТИНГИ ---
\lstset{
    basicstyle=\ttfamily\footnotesize,
    breaklines=true,
    texcl=true,
    frame=none
}